\documentclass[11pt, letterpaper]{article}

%% ─── PACKAGES ───────────────────────────────────────────────────────────────
\usepackage[top=1in, bottom=1in, left=1in, right=1in]{geometry}
\usepackage{setspace}          % \doublespacing
\usepackage{titlesec}          % numbered section headings
\usepackage{fancyhdr}          % running header / footer
\usepackage{microtype}         % microtypographic refinements
\usepackage{enumitem}          % fine-grained list control
\usepackage[hidelinks]{hyperref}

%% ─── PARAGRAPH & SPACING ────────────────────────────────────────────────────
\setlength{\parindent}{0.5in}
\setlength{\parskip}{0pt}

%% ─── SECTION HEADING FORMAT (numbered, bold) ────────────────────────────────
\titleformat{\section}{\normalfont\large\bfseries}{\thesection.}{0.5em}{}
\titlespacing*{\section}{0pt}{18pt}{6pt}

%% ─── LIST DEFAULTS ──────────────────────────────────────────────────────────
\setlist{nosep, itemsep=4pt, parsep=0pt, topsep=6pt}

%% ─── HEADER / FOOTER ────────────────────────────────────────────────────────
\pagestyle{fancy}
\fancyhf{}
\lhead{\small Kamrul, Md Kamruzzaman}
\rhead{\small EDCI 59100-016 $|$ Summer 2026}
\cfoot{\thepage}
\renewcommand{\headrulewidth}{0.4pt}
\setlength{\headheight}{14pt}

%% ════════════════════════════════════════════════════════════════════════════
\begin{document}
%% ════════════════════════════════════════════════════════════════════════════

%% ─── TITLE BLOCK ────────────────────────────────────────────────────────────
\begin{center}
  {\Large\bfseries Research Topic and Methodology Justification}\\[0.75em]
  {\normalsize EDCI 59100-016 $|$ Summer 2026}\\[0.25em]
  {\normalsize Md Kamruzzaman Kamrul}\\[0.25em]
  {\normalsize June 3, 2026}
\end{center}

\bigskip
\doublespacing

%% ─────────────────────────────────────────────────────────────────────────────
\section{Research Topic}
%% ─────────────────────────────────────────────────────────────────────────────

\noindent\textbf{Topic:} Using Vision-Language Models to support spatiotemporal
reasoning for safety monitoring and activity tracking on active construction sites.

Construction sites rank among the most hazardous occupational environments
globally. They are characterized by continuous, often unpredictable interactions
among workers, heavy equipment, and materials across a spatially and temporally
evolving environment. Incidents arising from undetected unsafe behaviors, such
as personal protective equipment (PPE) non-compliance, unauthorized proximity
to machinery, or deviations in task sequences, carry severe human and economic
consequences. At the same time, accurate progress tracking remains a persistent
operational challenge: manual inspection and reporting introduce delays and
inconsistencies that compound project delivery risk. This review examines how
Vision-Language Models (VLMs) and related multimodal architectures have been
applied to these two intertwined problems within active construction environments
and their simulated or laboratory-reconstructed equivalents. The topic sits at
the intersection of construction management, computer vision, and natural
language processing, making it relevant to researchers and practitioners across
all three communities.

%% ─────────────────────────────────────────────────────────────────────────────
\section{Research Questions and Purpose}
%% ─────────────────────────────────────────────────────────────────────────────

The purpose of this review is to systematically map the current state of
Vision-Language Model research as applied to spatiotemporal reasoning in
construction---identifying prevalent architectures, data fusion strategies,
deployment barriers, and evaluation frameworks, and locating the gaps where
future primary research is most needed. To operationalize this purpose, the
review is structured around four specific research questions (RQs):

\begin{itemize}[leftmargin=0.75in, label={}]
  \item \textbf{RQ1:} What are the prevalent VLM architectures currently
        utilized for spatiotemporal analysis in construction environments?
  \item \textbf{RQ2:} How is multimodal data---visual and textual---fused and
        processed to evaluate safety hazards on active construction sites?
  \item \textbf{RQ3:} What are the primary methodological bottlenecks in
        deploying these models for real-time progress tracking and safety
        monitoring?
  \item \textbf{RQ4:} How is the performance of VLMs currently evaluated and
        benchmarked for spatiotemporal tasks in construction contexts?
\end{itemize}

\noindent Together, these four questions guide a structured mapping of methods,
datasets, evaluation practices, and deployment realities across the existing
literature.

%% ─────────────────────────────────────────────────────────────────────────────
\section{Selected Review Methodology}
%% ─────────────────────────────────────────────────────────────────────────────

\noindent\textbf{Methodology:} Scoping review, conducted in accordance with
PRISMA-ScR (Preferred Reporting Items for Systematic Reviews and
Meta-Analyses extension for Scoping Reviews) guidelines.

A scoping review is a form of knowledge synthesis that systematically locates
and charts the available evidence across a defined topic area. Rather than
aggregating effect sizes or appraising the methodological quality of individual
studies---as a systematic review or meta-analysis would do---a scoping review
prioritizes breadth, comprehensive coverage, and transparent, reproducible
reporting. The goal is to characterize the range and nature of the existing
evidence, identify where the literature is sparse or inconsistent, and clarify
how key concepts and methods are being used across disciplines. The PRISMA-ScR
framework provides a structured checklist and flow diagram that makes every
stage of the search, screening, and selection process fully auditable.

%% ─────────────────────────────────────────────────────────────────────────────
\section{Methodology Justification}
%% ─────────────────────────────────────────────────────────────────────────────

The scoping review methodology is the most defensible and appropriate choice
for this topic and research questions for two principal reasons, in addition
to several practical considerations.

\textbf{Alignment with mapping-oriented research questions.} All four research
questions ask ``what'' and ``how''---they seek a structured map of
architectures, fusion strategies, bottlenecks, and evaluation frameworks rather
than a quantitative estimate of a particular intervention effect. Scoping
reviews are purpose-built for this kind of inquiry. A traditional meta-analysis
would be premature at this stage because the primary objective here is to chart
the extent, range, and nature of an emerging architectural family, not to
synthesize narrow intervention effects across comparable controlled experiments.

\textbf{Accommodation of methodological heterogeneity.} Evaluation metrics
across the relevant literature vary substantially: from geometric precision
metrics such as mean Average Precision (mAP), used in object-detection studies,
to natural language generation metrics such as BLEU and CIDEr, used in image
captioning and visual question-answering studies. This heterogeneity renders a
standardized statistical meta-analysis inappropriate. A scoping review
explicitly accommodates diverse study designs and reporting conventions within
one coherent synthesis framework, without requiring methodological equivalence
as a precondition for inclusion.

\textbf{Breadth of literature included.} Because the topic spans construction
management, computer vision, and natural language processing, relevant studies
include empirical field deployments, simulation-based benchmarks, dataset
introduction papers, and foundational architectural comparisons. The scoping
methodology permits inclusion of all these study types, which is essential for
capturing the full landscape of a multidisciplinary and fast-moving field.

\textbf{Practical considerations.} The PRISMA-ScR framework makes the review
transparent and reproducible within the constraints of this course. The
structured search strings, two-stage screening procedure, and flow diagram
documentation can be executed using the databases available through the
university library, within the available timeframe, and with results that
are fully auditable by an external reader.

%% ─────────────────────────────────────────────────────────────────────────────
\section{Reflection on the Search Process}
%% ─────────────────────────────────────────────────────────────────────────────

Three significant developments have shaped this project since the initial
proposal. First, in response to instructor feedback from Week~1, the original
broad framing---VLMs in built environments generally---was narrowed to focus
specifically on safety monitoring and activity tracking during the active
construction phase. This narrowing made the research questions more precise
and the inclusion criteria more coherent and defensible. Second, a single
broad research question was replaced with four operationally distinct questions
(RQ1--RQ4) that distribute the mapping task across architectures, fusion
strategies, deployment barriers, and evaluation frameworks. This structure
makes the eventual synthesis more organized and the contribution of the review
clearer to a reader. Third, the adoption of the PRISMA-ScR reporting framework
emerged as a natural commitment once the scoping methodology was confirmed,
because it provides the reproducibility scaffolding that any peer-reviewable
evidence synthesis requires and formalizes what would otherwise be informal
methodological choices.

%% ════════════════════════════════════════════════════════════════════════════
\end{document}
%% ════════════════════════════════════════════════════════════════════════════